% =============================================================================
% §4 EXPERIMENTAL PROTOCOL
% Every choice below is preregistered (docs/preregistration.md, §2/§3 + §6
% amendments). Nothing here may be decided after seeing a number; if a choice
% must change, amend the preregistration first (append-only) and say so here.
%
% Evaluation is a REFUTATION TEST (playbook §6.3): each experiment is designed
% so that, if our claim is wrong, the experiment shows it.
% =============================================================================
\section{Experimental Protocol}
\label{sec:protocol}

\subsection{Preregistration}

We call the protocol by its honest name: a confirmatory analysis protocol for
exploratory findings. The degeneracy and surrogate-metric observations were first
noticed on an earlier pipeline that carried a budget defect, a defective
baseline, and self-rolled metrics; every number from that pipeline is discarded
and none is reused here. What remains from it is a set of hypotheses, and the
protocol below pre-specifies their confirmatory re-analysis: the decision metric,
success threshold, and analysis for each claim were fixed, in an append-only
document, before the clean grid was run, so that no choice in this section could
be made after seeing a number. \TODO{author task at submission: cite the public
OSF/Zenodo registration DOI of that protocol document.}

\subsection{Data and Cohort}

We use the BraTS~2020 training collection, the only split with public
voxel-level ground truth. Each patient contributes a single two-dimensional
slice, chosen a priori as the slice of largest whole-tumor area, so that the
input given to the thresholding methods and to the learned baseline is identical.
The cohort is every patient for which that slice is defined: $368$ of the $369$
released cases, one case being excluded by the a priori tumor-slice criterion.
The cohort manifest, with a per-case checksum, is committed with the code.

% ⚠️ Never claim "state of the art on the BraTS leaderboard": the real test set
% is held out on Synapse. We use the training split and split it ourselves.
% ⚠️ A8: the LGG cohort may NOT be external — BraTS 2020 absorbed TCGA-LGG cases
% from the same collection. name_mapping.csv MUST be cross-checked and the
% overlap REPORTED AS A NUMBER before the words "external validation" are used.
We evaluate on the publicly annotated training split and construct our own
patient-level splits; we make no claim relating to the withheld BraTS test set.
Two targets are studied. The primary is whole tumor on FLAIR, the most favorable
case for the methods under audit: a single bright region that a contiguous
intensity band can plausibly capture. The hard target is enhancing tumor on
T1ce, a bright rim around a dark necrotic core, a shape no single contiguous
intensity band can represent. The argument is a fortiori: a limitation that
appears even on the favorable target is not an artifact of a poorly chosen case,
and the hard target shows where the representational limit is provable rather
than measured.

\subsection{Splits and Leakage Control}

% 🔴 A3: this subsection exists because the positive contribution was, as
% originally planned, train-on-test. Inside a paper that criticizes unsound
% comparison, that is the ONE error that gets rejected for INTEGRITY.
Splitting is at the \textbf{patient} level, never the slice level.
\nfolds-fold outer cross-validation, stratified by grade and tumor-volume
tertile, with an inner validation fold inside each outer-training fold for
epoch selection. Methods are classified once, a priori, into three families and
evaluated accordingly:

\begin{itemize}
  \item \textbf{Family A --- unsupervised per-image methods}: the classical
  thresholds (Otsu, Li, triangle), $k$-means and GMM clustering, the exact
  dynamic program, the metaheuristics, and random search. These have no fitted
  parameters and are evaluated on every patient.
  \item \textbf{Family B --- methods with learned parameters}: the 2-D U-Net,
  the single-parameter percentile threshold, and any data-driven selection of the
  threshold count $k$. These are evaluated \emph{out-of-fold only}; scoring them
  on patients used to fit them would be the exact error this paper documents.
  \item \textbf{Family C --- oracles}: the contiguous-band and gray-level-set
  oracles. These are not methods but unattainable upper bounds; they consume the
  test-time ground truth and are labeled ``uses test-time ground truth'' in every
  table and figure.
\end{itemize}

Nested cross-validation yields exactly one out-of-fold value per patient for
every method, on the same patient set, which is what makes paired Wilcoxon,
Friedman, and equivalence testing well defined at all. Intensity normalization
is computed per image; dataset-level statistics would leak through preprocessing.
The ``identical input'' constraint on the learned baseline is a constraint at
\emph{inference}: the network is trained on the tumor-bearing slices of the
outer-training patients and evaluated on the single designated slice of each
held-out patient.

\subsection{Budget}

All metaheuristics, including the quantum-inspired variant and random search,
receive an identical hard budget of objective function evaluations. Every
evaluation passes through a single gate, so that opposition-based
initialization, Lévy steps and memetic refinement are counted; an assertion
requires each algorithm to consume exactly the same number of evaluations,
and the grid is void if it fails.
% This is the module that killed the earlier batch (a 13.4% budget surplus).
% tests/test_nfe_budget.py is a HARD GATE: fail => the grid must not run.
% Equal TUNING budget too, not just equal NFE (prereg §6/A6): apparent advantage
% in this field almost always comes from uneven tuning.

% A6: breaking the circularity. We must not label an implementation "broken"
% because it fails to reach the optimum — that is the very hypothesis under test.
Implementation correctness is established \emph{independently} of the
thresholding hypothesis: every implementation must first reproduce published
multilevel-threshold benchmark tables on standard gray-level images at the
published evaluation budget. Exclusion is by this a priori criterion only, and
excluded implementations are reported separately rather than dropped.

\subsection{Metrics}

The clinical endpoints are Dice and the 95th-percentile Hausdorff distance,
with the Normalized Surface Dice at tolerance $\nsdtau$~mm as primary and a
sensitivity analysis over $\{1,3,5\}$~mm; overlap and surface metrics use a named
public implementation, since a paper auditing evaluation may not rely on an
unnamed or self-rolled one. PSNR and SSIM are computed from a standard library
and reported as evidence about the surrogate criteria, never as quality results.
We also report the connected-component count of the predicted mask and the
empty-mask rate as an independent result. IoU is omitted as a monotone bijection
of Dice. Because Dice is bounded and skewed, it is summarized by its median with
interquartile range and a bootstrap confidence interval, not by mean and standard
deviation, and the per-patient distribution is provided.

% A7 conventions, locked before running:
% - empty predicted mask with non-empty GT => Dice 0, HD95 = declared fixed
%   penalty (image diagonal). Both empty => Dice 1, HD95 0.
% - 🔴 It is FORBIDDEN to silently drop empty-mask cases from HD95 statistics:
%   that is selection bias favoring thresholding exactly at large k.
% - mean±std is the WRONG summary for Dice (bounded, skewed, often bimodal):
%   report median [IQR] + bootstrap CI + the per-patient distribution.

\subsection{Statistics}

Paired comparisons use the Wilcoxon signed-rank test at the patient level with a
rank-biserial effect size; omnibus ranking uses Friedman with a Nemenyi
critical-difference diagram~\cite{demsar_cd}. Equivalence is tested by two
one-sided tests and reported as a $90\%$ confidence interval of the paired
difference together with the smallest margin at which equivalence holds, so that
the analysis returns a number rather than a pass/fail verdict; a Bayesian
signed-rank test with a region of practical equivalence is the primary
equivalence analysis~\cite{benavoli_bayes}. We state a SESOI hierarchy
($\sesoistrict$ strict, an inter-rater-anchored intermediate
bound~\cite{menze_brats}, and $\sesoilenient$ lenient) and report all three.
Multiple-testing families are declared explicitly; the intersection--union
equivalence tests are not corrected. The primary threshold count is $\kprimary$,
the primary decoding rule is band selection of the brightest class, and the
reference for equivalence is the exact dynamic program.

% 🔴 A4d — THE scipy TRAP. Under the degeneracy mechanism, two algorithms that
% find the same threshold produce bit-identical masks => paired difference is
% EXACTLY 0. scipy.stats.wilcoxon defaults to zero_method='wilcox', which DROPS
% all zero pairs: effective n collapses, and the test silently becomes a
% statement about the minority of patients where the algorithms disagree — the
% opposite of the paper's message, and a manufactured "significance" of exactly
% the kind this paper documents.
% LOCKED: zero_method='pratt', and report n_zero/n_total in EVERY paired table.
Tied pairs are handled by Pratt's method rather than being discarded, and the
number of exactly tied patients is reported in every paired comparison; under
the mechanism of §\ref{sec:degeneracy} such ties are the expected outcome, not a
nuisance.

% --- FIGURE 1 — pipeline + budget diagram (drawio-skill, manual) -------------
% Provenance: docs/RESULTS.md "Hình 1". Flow figures are drawn, not bitmapped
% (playbook §6.4).
\begin{figure}[t]
  \centering
  \begin{tikzpicture}[
    font=\footnotesize,
    node distance=6mm and 7mm,
    box/.style={draw, rounded corners=2pt, align=center, inner sep=3pt,
                minimum height=8mm, text width=17mm},
    stage/.style={box, fill=blue!6},
    oracle/.style={box, fill=gray!12},
    score/.style={align=center, font=\scriptsize\itshape, text=black!65},
    arr/.style={-{Latex[length=2mm]}, thick}]
    \node[box] (hist) {Image\\ histogram $h$};
    \node[stage, right=of hist] (opt) {Stage 1\\ optimizer $\mathcal{A}$};
    \node[box, right=of opt] (thr) {thresholds\\ $\hat{\mathbf t}$};
    \node[stage, right=of thr] (dec) {Stage 2\\ decoding $D$};
    \node[box, right=of dec] (mask) {binary\\ mask $M$};
    \node[oracle, below=9mm of thr] (gt) {reference\\ mask $G$};
    \draw[arr] (hist) -- (opt);
    \draw[arr] (opt) -- (thr);
    \draw[arr] (thr) -- (dec);
    \draw[arr] (dec) -- (mask);
    \draw[arr] (gt) -- (dec.south);
    \draw[arr] (gt.east) -- ++(2.2,0) |- (mask.east)
          node[pos=0.25, below, score]{Dice$(M,G)$};
    \node[score, above=1mm of opt] {scored by $F_{\mathrm{Kapur}}$};
    \node[score, above=1mm of dec] {selects the object band};
    % budget gate on the optimizer
    \node[draw, dashed, fit=(opt), inner sep=5pt, label={[font=\scriptsize]below:%
      equal NFE budget}] {};
  \end{tikzpicture}
  \caption{The two-stage pipeline. Stage~1 optimizes the threshold vector on the
  histogram and is scored in fitness units ($F_{\mathrm{Kapur}}$); Stage~2 decodes
  the intensity bands into a binary mask and is the only stage scored against the
  reference mask (Dice). A decade of work optimizes Stage~1. Every optimizer
  passes through a single evaluation gate enforcing an identical function-%
  evaluation budget. The reference mask enters only at Stage~2 and at scoring;
  fitting of any learned component is confined to the training split
  (§\ref{sec:protocol}).}
  \label{fig:pipeline}
\end{figure}
